\documentclass[12pt,addpoints,answers]{exam}
\usepackage[utf8]{inputenc}
\usepackage{amsmath,amsfonts,amssymb,amsthm}
\usepackage[margin=1in]{geometry}
\usepackage{mathtools}
\usepackage{hyperref}
\usepackage{fullpage}
\usepackage{microtype}
\usepackage{xspace}
\usepackage[svgnames]{xcolor}
\usepackage[sc]{mathpazo}
\usepackage{enumitem}
\usepackage{bm}

%----------Header--------------------%
\def\course{{\sc Introduccion a la Criptografia Moderna}}
\def\term{Depto. de Computaci\'{o}n, FCEN, UBA, 1er. Bimestre 2026}
\def\prof{Lecturer: Juan Garay}
\newcommand{\handout}[5]{
   \renewcommand{\thepage}{\arabic{page}}
   \begin{center}
   \framebox{
      \vbox{
    \hbox to 5.78in { \hfill \large{\course} \hfill }
    \vspace{2mm}
%    \hbox to 5.78in { \hfill \large{\prof} \hfill }
%       \vspace{2mm}
       \hbox to 5.78in { {\Large \hfill \textbf{#5}  \hfill} }
       \vspace{2mm}
       \hbox to 5.78in { \term \hfill \emph{#2}}
       \hbox to 5.78in { {#3 \hfill \emph{#4}}}
      }
   }
   \end{center}
   \vspace*{4mm}
}
\newcommand{\hw}[4]{\handout{#1}{#2}{#3}{#4}{Homework #1}}

% -- For ignoring stuff -- %
\newcommand{\ignore}[1]{}


\begin{document}

%----Specs: change accordingly-----%
\newif\ifstudent % comment out false
\studenttrue 
% \studentfalse

\def\hwnum{2}
\def\issuedate{14/4/26}
\def\duedate{23/4/25, 23:59 hs} % 
\def\yourname{} % put your name here
%------------------------------%

\ifstudent
\hw{\hwnum}{\issuedate}{Student: \yourname}{Due: \duedate}%
\else
\hw{\hwnum}{\issuedate}{\prof}{Due: \duedate}%
\fi

\noindent \textbf{Instructions}

\begin{itemize}
    \item Upload your solution to Campus; make sure it's only one file, and clearly write your name on the first page. Name the file \textsf{`$<$your last name$>$\_HW2.pdf.'} 
    %{\bf Important:} Make sure to tap {\bf Turn in} after you upload your solution.   
      
     If you are proficient with \LaTeX, you may also typeset your submission and submit in PDF format. To do so, uncomment the ``\%\textbackslash begin\{solution\}'' and ``\%\textbackslash end\{solution\}'' lines and write your solution between those two command lines.
    
      \item Your solutions will be graded on \emph{correctness} and
    \emph{clarity}. You should only submit work that you believe to be
    correct.
    % , and you will get significantly more partial credit if you     clearly identify the gap(s) in your solution.
    
    \item You may collaborate with others on this problem set.  However,
    you must \textbf{{write up your own solutions}} and \textbf{{list
      your collaborators and any external sources (including ChatGPT and similar generative AI chatbots)}} for each
    problem. Be ready to explain your solutions orally to a member of the course's staff
    if asked.
    
    \ignore{
    \item For problems that require you to provide an algorithm, you must
    give a precise description of the algorithm, together with a proof
    of correctness and an analysis of its running time. You may use
    algorithms from class as subroutines. You may also use any facts
    that we proved in class or from the book.
    } %IGNORE
    
\end{itemize}


\noindent This assignment contains \numquestions\ questions,
% \numpages\ pages
for a total of \numpoints \ points.
% and \numbonuspoints\ bonus points.

\bigskip
%\newpage

\begin{questions}


\question Let $F$ be a pseudorandom permutation. For each of the following encryption schemes (where the key is a uniform $k\in\{0,1\}^n$), show how to decrypt and then state whether the scheme is CPA-secure. Explain your answers.
\begin{parts}
    \part[5] To encrypt $m\in\{0,1\}^n$, choose uniform $r\in\{0,1\}^n$ and output the ciphertext $\langle r,F_k(m)\oplus r\rangle$.
    \begin{solution}
      Para decriptar, tenemos que reconocer los dos elementos de la tupla (asumimos una codificación estandar para esto). Luego, podemos hacer la inversa del XOR con $r$ (el primer elemento de la tupla) para obtener $F_k(m)$. Luego, como tenemos $k$, aplicamos la inversa de $F_k$ para obtener $m$.

      Vemos que este esquema no es CPA-seguro. Para demostrarlo mostramos una estrategia que puede seguir el adversario en el experimento de $Priv_k$. 
  
      El adversario, una vez que entrega los $m_0$ y $m_1$ y recibe la versión encriptada de uno de los dos, pide encriptar al oráculo $m_0$. Luego, a ambos ciphertexts los manipula para llegar a $F_k(m_0)$ y $F_k(m_b)$ (lo hace con el primer paso de la decriptación mostrada antes, es trivial para cualquier tupla). De ese modo, como $F_k$ es una función determinística, $F_k(m_0) = F_k(m_b)$ si $b = 0$, y si no son iguales es porque $b=1$. 
      Luego, el adversario tiene probabilidad mayor a $\frac{1}{2}$ de ganar el experimento y el esquema no es CPA-seguro.

    \end{solution}
    \part[5] To encrypt $m\in\{0,1\}^n$, choose uniform $r\in\{0,1\}^n$ and output the ciphertext $\langle r,F_r(k)\oplus m\rangle$.
    \begin{solution}
      Para decriptar, tenemos que reconocer los dos elementos de la tupla, $r$ y $F_r(k)\oplus m$. Luego, podemos generar $F_r(k)$ (dado que tenemos $k$) y hacer la inversa del XOR para obtener $m$. 

      Vemos que este esquema tampoco es CPA-seguro. Para demostrarlo veamos el siguiente escenario. El adversario entrega $m_0$ y $m_1$ y recibe $ c = \langle r, F_r(k) \oplus m_b \rangle $. Luego, conoce $r$, y $F$ es una función pública. Pide al oráculo que encripte un $m_2$ tal que $m_2 = 0^n$ y recibe $ c' = \langle r_2, F_{r_2}(k) \oplus 0^n \rangle $. Entonces obtiene $F_{r_2}(k)$ haciendo el XOR con $0^n$.
      
      Luego, para obtener $k$ puede hacer lo siguiente. Como $F$ es pública, tiene $r_2$ y tiene $F_{r_2}(k)$, $k = F_{r_2}^{-1}(F_{r_2}(k))$. 
     
      Luego, volviendo al ciphertext original, $c = \langle r, F_r(k) \oplus m_b \rangle $, puede obtener $m_b$ haciendo el XOR con $F_r(k)$, dado que ahora conoce $k$, $r$ y $F$. Entonces compara $m_b$ con $m_0$ y $m_1$ y gana el experimento con probabilidad 1.
      
    \end{solution}
\end{parts}


\question[10] Show that the encryption scheme presented in Lecture 5 (p. 21) is CCA-secure.
\begin{solution}
  La diferencia entre CCA-seguro y CPA-seguro es que en CCA-seguro el adversario puede decriptar cualquier ciphertext (excepto el target ciphertext).
  
  En este caso, vemos que el adversario recibe un $c^*$. Las opciones que tiene son: 
  \begin{itemize}
    \item Encriptar mensajes arbitrarios.
    \item Decriptar ciphertexts arbitrarios (posibles modificaciones de $c^*$, o ciphers de mensajes arbitrarios). 
  \end{itemize}

  Si encripta mensajes arbitrarios, va a recibir un ciphertext. Desencriptarlo sin modificaciones va a devolver el mismo mensaje, lo que no aporta nueva información. Mas formalmente, si toma $m'$ y pide encriptar tal que $c' = Enc_k(m')$, pedir decriptar $c'$ va a dar $m' = Dec_k(c')$. 

  El cipher en sí es el resultado de una PRP, en el cual la clave es secreta. Entonces no hay mucho que se pueda inferir sobre el contenido de $x$, $r$ o $0^n$ a partir de un ciphertext arbitrario. 

  Si el adversario intenta modificar $c^*$ para luego intentar decriptarlo y obtener información, tiene probabilidades negligibles de encontrar un cipher válido. En particular, de $\frac{1}{2^{n}}$, dado que los ciphers válidos son los que terminan en $0^n$. Además, si encuentra un cipher que sea válido, con mínimas modificaciones a $c^*$, no quiere decir que ese cipher modificado tenga el mismo mensaje $x$. 
  
  Por ende, vemos que agregar el oráculo de decripción no aporta nada. Con ver que es CPA-seguro alcanza. Y esto es bastante simple de ver, ya que $F_k$ es PRP, y lo que se encripta recibe siempre un $r$ fresco uniforme de input que hace que las encripciones sean indistinguibles unas de otras. 
\end{solution}


\question Prove that the following modifications of basic CBC-MAC {\bf do not} yield a secure MAC (even for fixed-length messages):
\begin{parts}
    \part[5] {\sf Mac} outputs all blocks $t_1,\ldots,t_\ell$, rather than just $t_\ell$. (Verification only checks whether $t_\ell$ is correct.)
    \begin{solution}
      Para mostrar que el esquema no es seguro, mostramos una estrategia que permite ganar el experimento.

      Sea $m = m_1 || m_2 || ... || m_l$. Pido al oráculo $MAC_k(m) = (t_1, t_2, ... , t_l)$. Hago lo mismo para $m' = m'_1 || m_2 || ... || m_l$, y pido al oráculo $MAC_k(m') = (t_1', ..., t'_l)$. Notar que $t_1 = F_k(m_1)$ y $t'_1 = F_k(m'_1)$.

      Ahora definimos otro mensaje $m^* = m'_1 || t_1' \oplus t_1 \oplus m_2 || m_3 || ... || m_l$. 

      Luego, vamos a construir el $MAC_k(m^*)$ sin pedirlo al oráculo de esta manera. 

      $$t^*_1 = F_k(m'_1) = t'_1$$
      $$t^*_2 = F_k(t'_1 \oplus t'_1 \oplus t_1 \oplus m_2) = F_k(t_1 \oplus m_2) = t_2$$
      $$t^*_3 = F_k(t_2 \oplus m_3) = t_3$$

      Y así sucesivamente hasta completar los $l$ bloques. De ese modo, construimos el tag de un mensaje $m^*$ sin haberlo pedido al oráculo. 
    \end{solution}
    \part[5] A random initial block is used each time a message is authenticated. That is, choose uniform $t_0\in\{0,1\}^n$, run basic CBC-MAC over the ``message'' $t_0,m_1,\ldots,m_\ell$, and output the tag $\langle t_0,t_\ell\rangle$. Verification is done in the natural way.
    \begin{solution}

      Tomamos un mensaje $m = m_1 || ... || m_l$. Luego pedimos su MAC y conseguimos $\langle t_0, t_l \rangle$.

      Notamos que $t_1 = F_k(t_0 \oplus m_1)$ por la construcción del esquema, donde $t_0$ fue elegido al azar de manera uniforme por el oráculo de MAC.

      Luego, definimos un mensaje $m^* = m_1^* || m_2 || m_3 || ... || m_l$. Le modificamos solo el primer bloque para ver si le podemos generar un tag sin pedirle al oráculo, solamente con la información de $m$. Para que sean distintos, necesito $m_1^* \neq m_1$

      Veamos que necesitamos pedir lo siguiente:
      $$t_1^* = F_k(t_0^* \oplus m_1^*) = t_1 = F_k(t_0 \oplus m_1)$$
      Si se cumple esto, el resto de $t_i^*$ van a ser iguales a $t_i$ para cada $i=2...l$.

      Digamos que $m_1^* = x \oplus m_1$, para algún $x \in \{0, 1\}^n$, con $x \neq 0^n$ (así es distinto a $m^*$). Luego, tenemos que $$F_k(t_0^* \oplus x \oplus m_1) = F_k(t_0 \oplus m_1)$$ 
      Luego, para anular el $x$, podemos pedir que $t_0^* = t_0 \oplus x$.
      De ese modo, $t_1^* = F_k(t_0 \oplus x \oplus x \oplus m_1) = F_k(t_0 \oplus m_1) = t_1$.
      
      Luego, conseguimos $t_l$ aplicando el resto del algoritmo, y $t_0^*$ es $t_0 \oplus x$. De ese modo, podemos generar el tag de $m^*$ sin pedirle al oráculo, es decir que $Vrfy(\langle t_0^*, t_l^* \rangle, m^*) = 1$.

      
    \end{solution}
\end{parts}


\question
\begin{parts}

\part[5] Provide formal security definitions for {\em second preimage resistance} and {\em preimage resistance} of hash functions, similar to the definition of collision resistance we saw in Lecture 6. 
\begin{solution}
  Para definir resistencia a la segunda preimagen, definimos el siguiente experimento, similar al de resistencia a la colisión. En este experimento y el siguiente, la función de hash es $\mathsf{H} = (Gen, H)$.

  \begin{enumerate}
    \item Se genera una clave $s$ con $Gen(1^n)$
    \item El adversario $\mathsf{A}$ recibe $s$ y $x$, un mensaje elegido de manera uniforme por el challenger. Notamos que si $H$ es de longitud fija, $x$ va a tener la longitud correspondiente. Luego, el adversario devuelve $x'$.
    \item El output del experimento es $1$ $\iff$ $H^s(x') = H^s(x)$ y $x' \neq x$.
  \end{enumerate}

  Y definimos la resistencia a la segunda preimagen como:

  Una función $\mathsf{H} = (Gen, H)$ es resistente a la segunda preimagen si, para todo adversario PPT $\mathsf{A}$, existe una función negligible $negl(n)$ tal que:

  $$
    \Pr[\text{SecondPreimage}_{\mathsf{A}, \mathsf{H}}(n) = 1] \leq negl(n)
  $$


  Veamos ahora la resistencia a la preimagen. Definimos el siguiente experimento:
  \begin{enumerate}
    \item Se genera una clave $s$ con $Gen(1^n)$
    \item El adversario $\mathsf{A}$ recibe $s$ e $y$ tal que $y = H^s(x)$ (con $x$ elegido aleatoriamente de manera uniforme). Luego, el adversario devuelve $x'$.
    \item El output del experimento es $1$ $\iff$ $y = H^s(x')$.
  \end{enumerate}

  Y definimos la resistencia a la preimagen como:
  $$
    \Pr[\text{Preimage}_{\mathsf{A}, \mathsf{H}}(n) = 1] \leq negl(n)
  $$

\end{solution}
\part[5] Prove that any hash function that is collision resistant is second-preimage resistant, and any hash function that is second-preimage resistant is preimage resistant.
\begin{solution}
  Veamos primero la primera implicación. Asumimos que tenemos una función $H$ que es resistente a la colisión y veamos que cumple con la resistencia a la segunda preimagen.

  Vamos a demostrarlo por el absurdo: asumimos que H no es resistente a la segunda preimagen, es decir, que sí es posible encontrar, dado un $x$, un $x'$ tal que $H(x') = H(x)$ y $x' \neq x$. Pero vemos que esto significaría que pudimos encontrar un par $x, x'$ tal que $H(x) = H(x')$ (con $x \neq x'$), contradiciendo la resistencia a la colisión de $H$. Absurdo! Luego, H es resistente a la segunda preimagen.

  Vamos con la segunda implicación. Asumimos que $H : \{0,1\}^{l'} \to \{0,1\}^{l}$ comprime, con $l' = 2l$ (alcanza con $l' - l = \omega(\log n)$). Lo demostramos por el absurdo: supongamos que existe un adversario PPT $\mathsf{B}$ que gana el experimento de resistencia a la preimagen con probabilidad no negligible $\epsilon(n)$. Es decir, dados $s$ e $y = H^s(x)$ con $x$ uniforme, $\mathsf{B}(s, y)$ devuelve con probabilidad $\epsilon(n)$ algún $x'$ con $H^s(x') = y$ (no necesariamente $x' \neq x$).

  Construimos un adversario PPT $\mathsf{A}$ para el experimento de resistencia a la segunda preimagen:
  \begin{enumerate}
    \item Recibe $(s, x)$ con $x \leftarrow \{0,1\}^{l'}$ uniforme.
    \item Computa $y := H^s(x)$.
    \item Computa $x' := \mathsf{B}(s, y)$.
    \item Si $H^s(x') = y$ y $x' \neq x$, devuelve $x'$. Sino, $\bot$.
  \end{enumerate}

  Definimos los eventos (sobre la aleatoriedad de $x$, $s$ y $\mathsf{B}$):
  \begin{itemize}
    \item $W$: $\mathsf{B}$ tuvo éxito, es decir $H^s(x') = y$.
    \item $E$: $x \in S$, donde $S = \{x \in \{0,1\}^{l'} : |(H^s)^{-1}(H^s(x))| \geq 2\}$ son los inputs que \emph{tienen} una segunda preimagen.
  \end{itemize}
  Notar que $\mathsf{A}$ gana exactamente cuando ocurre $W \wedge E \wedge (x' \neq x)$ (si $x \notin S$, la única preimagen es $x$, y $W$ forzaría $x' = x$).

  Cota sobre $\Pr[E]$. El complemento $\overline{S}$ tiene a lo sumo una preimagen por cada $y \in \{0,1\}^l$, luego $|\overline{S}| \leq 2^l$ y
  $$\Pr[x \notin S] \leq \frac{2^l}{2^{l'}} = 2^{l-l'} = 2^{-l},$$
  que es negligible en $n$.

  Cota sobre $\Pr[x' \neq x \mid W, E]$. $x$ es uniforme sobre $\{0,1\}^{l'}$, y por lo tanto, condicionado a $s$ y a $y = H^s(x)$, $x$ queda distribuido uniformemente sobre $(H^s)^{-1}(y)$. La salida $x'$ de $\mathsf{B}$ depende sólo de $(s, y)$ (no de $x$). Entonces, fijados $s, y, x'$,
  $$\Pr[x' = x \mid s, y, x \in S] \leq \frac{1}{|(H^s)^{-1}(y)|} \leq \frac{1}{2}.$$
  Luego $\Pr[x' \neq x \mid W, E] \geq \frac{1}{2}$.

  Juntando todo:
  \begin{align*}
    \Pr[\mathsf{A} \text{ gana}] &= \Pr[W \wedge E \wedge (x' \neq x)] \\
    &= \Pr[W \wedge E] \cdot \Pr[x' \neq x \mid W, E] \\
    &\geq \tfrac{1}{2} \cdot \Pr[W \wedge E] \\
    &\geq \tfrac{1}{2} \cdot (\Pr[W] - \Pr[\overline{E}]) \\
    &\geq \tfrac{1}{2} \cdot (\epsilon(n) - 2^{-l}).
  \end{align*}
  Como $\epsilon(n)$ es no negligible y $2^{-l}$ es negligible, $\Pr[\mathsf{A} \text{ gana}]$ es no negligible. Absurdo con la hipótesis de que $H$ es resistente a la segunda preimagen. Luego, $H$ es resistente a la preimagen.

\end{solution}
\end{parts}

\question[10] Show how to find a collision in the Merkle tree construction if $t$ is not fixed. Specifically, show how to find two sets of inputs $x_1,\ldots,x_t$ and $x'_1,\ldots,x'_{2t}$ such that ${\cal MT}_t(x_1,\ldots,x_t)={\cal MT}_{2t}(x'_1,\ldots,x'_{2t})$.
\begin{solution}
  Para encontrar estos dos conjuntos, vamos a hacer lo siguiente. Vamos a "extender hacia abajo" el merkle tree chico, de $t$ hojas, para obtener un merkle tree grande, de $2t$ hojas. 

  Asumimos que tenemos un conjunto cualquiera de $x'_1...x'_{2t}$. Definimos entonces $x_i = H(x'_{2i-1}) || H(x'_{2i})$ para $i = 1...t$. De ese modo, la i-ésima hoja del $MT_t$ es $h_i = H(x_i) = H(H(x'_{2i-1}) || H(x'_{2i}))$, que es exactamente lo mismo que $h_{2i-1,2i}$ del $MT_{2t}$. Vemos entonces que si comparamos el $MT_t$ con el $MT_{2t}$ recortado en el último nivel (es decir, sacándole todas las hojas), son exactamente iguales! 
  Luego, van a tener la misma raíz. Luego, pudimos encontrar dos conjuntos $x_1...x_t$ y $x'_1...x'_{2t}$ tales que $MT_t(x_1...x_t) = MT_{2t}(x'_1...x'_{2t})$.
\end{solution}


\end{questions}


Nota: usé claude opus 4.7 para ayudarme con la demostración del ejercicio 4b, más que nada con la parte de las probabilidades de la última implicación. 

\end{document}
